\section{Tensor Product of OGS configurations}
\label{sec:tensor}

Following the introduction of the tensor product of configurations
for C-OGS in Definition~\ref{d:tensor_cogs}, we 
introduce a similar definition for the Alternating and Well-Bracketed LTS.
We also prove results relating the 
(complete) traces generated by $\FF \otimes \GG$
with the interleaving of traces generated by $\FF$ and $\GG$.


\subsection{Interleaving of complete traces for \COGS}

We first state the variant of Lemma~\ref{l:COGS-interleave-tensor}
for complete traces.
So given a complete trace for a configuration, the projections of the
trace on sub-configurations are also complete, as proven in the following lemma.


\begin{lemma}
  \label{l:COGS-decomp-ct}
  Suppose 
  \begin{itemize}
  \item $\FF = \appendCon{\FF_1}{\FF_2}$;
  \item $\FF \ArrC \tr \FF'$ and $ \tr $ is a complete trace
  for $\FF$.
  \end{itemize} 
  Then for both $i \in \{1,2\}$,
  there exists $\FF_i'$ and a complete trace $\tr_i$
  for $\FF_i$ such that
  \begin{enumerate}
  \item ${\FF_i} \ArrC{\tr_i} \FF_i'$;
  \item $\tr \in \interl{\tr_1}{\tr_2}$;
  \item $\FF' = \appendCon{\FF'_1} {\FF'_2}$.
  \end{enumerate}
\end{lemma}              
  
\begin{proof}
  From Lemma~\ref{l:COGS-interleave-tensor}(\ref{l:COGS-interleave-tensor_1}) 
  we get the existence of configurations $\FF_i'$ and traces $\tr_i$ 
  such that the three items above are verified.
  From Lemma~\ref{l:OGS_complete-strongly-passive}, we get that $\FF'$
  is strongly passive, so that both $\FF'_1,\FF'_2$ are strongly passive as well.
  Thus applying Lemma~\ref{l:OGS_complete-strongly-passive}
  in the other direction,
  we get that both $\tr_1,\tr_2$ are complete.
\end{proof} 

\begin{lemma}
  \label{l:COGS-recomp-ct}
  Suppose 
  \begin{enumerate}
  \item $\FF = \FF_1 \otimes \FF_2$;   
  \item $\FF_i \ArrC{\tr_i} \FF_i'$ where 
  $\tr_i$ is a complete trace for $\FF_i$, %\wrt $\stti$,
   $i=1,2$. 
  \end{enumerate} 
  Then there is a trace 
   $\tr \in \interl{\tr_1}{\tr_2}$ complete for $\FF$ such that
  $\FF \ArrC{\tr}{\FF'_1} \otimes {\FF'_2}$. 
\end{lemma}

\begin{proof}
  The existence of $\tr$ such that $\FF \ArrC{\tr}{\FF'_1} \otimes {\FF'_2}$ follows from 
  Lemma~\ref{l:COGS-interleave-tensor}(\ref{l:COGS-interleave-tensor_2}). 
  Since $\tr$ is the interleaving of two complete traces respectively from 
  $\FF_1$ and $\FF_2$, we get that it is complete for $\FF = \FF_1 \otimes \FF_2$.
\end{proof} 

\subsection{\AOGS}
\label{a:tensor-aogs}

\begin{definition}
Taking $\FF_1,\FF_2$ two compatible \AOGS configurations, with 
 $\gamma_i,\phi_i$ their respective environment and set of names,
$\appendCon{\FF_1} {\FF_2}$ is defined as:
\begin{itemize}
 \item $\conf{M,p,\gamma_1 \cdot \gamma_2,\phi_1 \cup \phi_2}$
 when one of $\FF_1,\FF_2$  is active, with $M,p$ its toplevel term and continuation name;
 \item $\conf{\gamma_1 \cdot \gamma_2,\phi_1 \cup \phi_2}$
 when $\FF_1,\FF_2$  are passive.
\end{itemize}
\end{definition}
\noindent
We can then state the adaptation of Lemma~\ref{l:COGS-interleave-tensor} to \AOGS.
We write $\interlA{\tr_1}{\tr_2}$ for the subset of $\interl{\tr_1}{\tr_2}$
formed by \emph{alternating} traces, with the
extra condition that if one of $\tr_1,\tr_2$
starts with an output, then the resulting interleaved trace must also start with an output.

\begin{lemma}
\label{c:dec_GS}
Suppose $\FF_1,\FF_2$ are compatible \AOGS configurations, and at most one of the two is active.
Then the set of traces generated by $\FF_1 \otimes \FF_2$
is equal to the set of interleavings $\interlA{\tr_1}{\tr_2}$,
with $\FF_1 \ArrA{\tr_1}$ and $\FF_2 \ArrA{\tr_2}$.
That is:
\begin{enumerate}
  \item\label{l:dec_GS_1} if $\appendCon{\FF_1}{\FF_2} \ArrA{\tr} \FF$, 
  then, for $i=1,2$,  there are traces $ \tr_i$
  such that $ {\FF_i} \ArrA{\tr_i} \FF'_i$, 
  and $\tr  \in \interlA{ \tr_1}{\tr_2}$, 
  and $\FF = \appendCon{\FF'_1}{\FF'_2}$;
  \item\label{l:dec_GS_2} conversely, if for $i=1,2$, we have
  $\FF_i \ArrA{ \tr_i} \FF'_i$,  
  and $\tr  \in \interlA {\tr_1}{\tr_2}$, 
  then $\appendCon{\FF_1}{\FF_2} \ArrA{\tr} 
  \appendCon{\FF'_1} {\FF'_2}$.
\end{enumerate}
\end{lemma}  
\noindent
The statement of (\ref{l:dec_GS_2}) above can be deduced from
the corresponding result for \COGS, namely
Lemma~\ref{l:COGS-interleave-tensor},
combined with the characterization
of traces generated by the \AOGS LTS 
as being the \emph{alternating} traces of \COGS LTS.
However, such a proof does not work for (\ref{l:dec_GS_1}),
since we cannot deduce directly that $\tr_1$ and $\tr_2$ are alternating.
To prove this direction,
we use the correspondence between \AOGS and $\ArrN{}$ the \outputpr LTS 
established in Corollary~\ref{c:traces}, 
combined with Lemma~\ref{l:dec_piI} below, whose proof uses  the following lemma.


\begin{lemma}
\label{l:dec_piI_aux}
Suppose $P| Q  \ArrN \mu P'|Q$ and 
$P  \Arr \mu P'$. Then also $P  \ArrN \mu P'$.
\end{lemma} 
\begin{proof}
If $P  \ArrN \mu P'$ were not possible, because of the constraint on input components in
the \opLTS with respect to the ordinary LTS, then also 
$P| Q  \ArrN \mu P'|Q$ would not be possible (if $P$ is not input reactive,  then 
$P|Q$ cannot be so either). 
\end{proof} 


\begin{lemma}
\label{l:dec_piI}
Suppose $\FF_1,\FF_2$ are compatible configurations, and at most one of the two is active.
\begin{enumerate}
\item If $\encoConV {\FF_1} |\encoConV {\FF_2} \ArrN \tr  P_1 | P_2  $, 
  then, for $i=1,2$,  there are traces $ \tr_i $
  such that $\encoConV {\FF_i} \ArrN { \tr_i} P_i $, 
  and $ \tr \in \interlA { \tr_1   }{ \tr_2 }$;
\item
Conversely, if for  $i=1,2$, we have
  $\encoConV {\FF_i} \ArrN { \tr_i} P_i$,
  $\tr \in \interlA { \tr_1   }{ \tr_2 }$, 
  then also 
  $\encoConV {\FF_1} |\encoConV {\FF_2} \ArrN \tr  P_1 | P_2$.
\end{enumerate} 
\end{lemma}  

\begin{proof}
Proof of (1): 
the assertion follows by reasoning as in Lemma~\ref{l:dec_pi}.
Processes 
$\encoConV {\FF_1}$ and  $\encoConV {\FF_2}$ cannot interact and, by
Lemma~\ref{l:ci_trans}, the property is invariant under transitions. 
Thus each action (visible
or silent) that contributes to the given trace $\tr$ is derived using either  
rule {\trans{parL}} or  rule {\trans{parR}} on the outermost parallel composition. 
The former case represents  an action (in the ordinary LTS) from the left process, the
latter case an action from the right process. In either case, the action can be lifted to 
the  \opLTS using Lemma~\ref{l:dec_piI_aux}. 
By iterating the reasoning on each action of the trace  $\tr$
we obtain traces $\tr_i$ ($i=1,2$) with  
 $\encoConV {\FF_i} \ArrN { \tr_i} P_i$. 
% We know 
% \begin{equation}
% \label{e:encoConV}
% \encoConV {\FF_1} |\encoConV {\FF_2} \ArrN \tr  P_1 | P_2  
% \end{equation} 
% Hence, by the definition of the LTS $\ArrN{}$, we also have 
%  $$\encoConV {\FF_1} |\encoConV {\FF_2} \ArrPI \tr  P_1 | P_2  .$$
% We now apply Lemma~\ref{l:dec_pi} to infer  
%  $\encoConV {\FF_i} \ArrPI { \tr_i} P_i$.
% It also holds 
%  $\encoConV {\FF_i} \ArrN { \tr_i} P_i$
% (the  observations in  $\tr_i$
%   were possible in \reff{e:encoConV},  with an extra parallel component, 
% which  may only
% reduce the possibility of transitions in the LTS $\mapsto$). 
% \gj{We may detail the proof that $\encoConV {\FF_i} \ArrN { \tr_i} P_i$.}


For (2): From $\encoConV {\FF_i} \ArrN { \tr_i} P_i$, we deduce that
$\encoConV {\FF_i} \ArrPI{ \tr_i} P_i$.
So applying Lemma~\ref{l:dec_pi}, we get that 
$\encoConV {\FF_1} |\encoConV {\FF_2} \ArrPI \tr  P_1 | P_2$.
Moreover, from $\encoConV {\FF_i} \ArrN { \tr_i} P_i$ we also deduce
that each $\tr_i$ is alternating. Since at most one of the two
$\FF_1,\FF_2$ is active, we deduce that at most one of $s_1,s_2$ 
is $P$-starting. We have that $\tr$ is alternating since
$\tr \in \interlA{ \tr_1 }{ \tr_2 }$, and 
active iff one of the composing traces is $P$-starting. This means that we can assume 
that when an input transition of $\tr$ is performed the source process is 
input reactive.
\end{proof} 

%\fi
\noindent
We now state the corresponding lemmas for complete traces,
whose proofs are similar to Lemma~\ref{l:COGS-decomp-ct} and Lemma~\ref{l:COGS-recomp-ct}.
\begin{lemma}
\label{l:AOGS-decomp-ct}
Suppose
\begin{itemize}
  \item $\FF  = \appendCon{\FF_1}{\FF_2}$;
  \item $\tr$ is a complete trace for $\FF$,
  such that $\FF \ArrA \tr \FF'$.
\end{itemize} 
Then for both $i \in \{1,2\}$,
there exists $\FF_i'$ and a complete trace $\tr_i$
for $\FF_i$ such that 
\begin{enumerate}
  \item ${\FF_i} \ArrA{\tr_i} \FF_i'$;
  \item $\tr \in \interlA{\tr_1}{\tr_2}$; 
  \item $\FF' = \appendCon{\FF'_1}{\FF'_2}$.
\end{enumerate}
\end{lemma}              

\begin{lemma}
  \label{l:AOGS-recomp-ct}
  Suppose 
  \begin{enumerate}
    \item $\FF = \FF_1 \otimes \FF_2$;   
    \item $\tr_i$ a complete trace for $\FF_i$,
    such that $\FF_i \ArrA{\tr_i} \FF_i'$,
    $i=1,2$. 
  \end{enumerate} 
  Then there is 
   $\tr \in \interlA{\tr_1}{\tr_2}$ complete for $\FF$ such that
  $\FF \ArrA{\tr}{\FF'_1} \otimes {\FF'_2}$. 
\end{lemma}
\noindent
Next, we show in Lemma~\ref{l:decACN} a decomposition property for 
complete trace equivalence $\TEc$ \wrt $\otimes$ on configurations.
%This lemma also holds for $\TEpc$.
%
To prove it, we rely on the following lemmas
that show that the action of a configuration is 
determined by a specific singleton component, 
and it is the action itself that allows us to determine such a component.

\begin{lemma}
\label{l:decNsig}
Suppose $\FF \arrA\mu \FF'$.
Then there is a singleton $\FF_1$, and some $\FF_1',\FF_2,\FF_2'$ 
such that:
\begin{enumerate}
\item $\FF =
 {\FF_1} \otimes {\FF_2}$;
\item $\FF_1 \arrA\mu \FF_1'$;
\item $\FF' =
 {\FF_1'} \otimes {\FF_2}$.
\end{enumerate}
\noindent
Moreover, the P-\Supports and the structure of continuations of $\FF_1$ and $\FF_1'$  
are determined by 
$\mu$; that is, 
if  $\okC \FF\GG$ and
 $\GG \arrA\mu \GG'$, then for  
 the singleton $\GG_1$, and the configurations $\GG_1',\GG_2,\GG_2'$ such that
 $\GG = \appendCon{\GG_1}{\GG_2}$,
 $\GG_1 \arrA\mu \GG_1'$, and 
 $\GG' = \appendCon{\GG_1'}{\GG_2}$ we have: 
 $\okC{\FF_1}{\GG_1}$, $\okC{\FF_2}{\GG_2}$, and 
$\okC{\FF_1'}{\GG_1'}$.  
 \end{lemma} 

\begin{proof}
By a simple case analysis on the form of $\mu$.
\end{proof} 
\noindent
The next lemma uses the previous result to show that the decomposition of a trace 
is uniform.

\begin{lemma}
\label{l:decNsig2}
Suppose 
\begin{enumerate}
\item
$\FF \ArrA \tr   \FF'$
and $\GG \ArrA \tr   \GG'$;
\item  
 $\FF = \FF_1 \otimes \FF_2$ and 
 $\GG = \GG_1 \otimes \GG_2$;  
\item 
  $\okC{\FF}{\GG}$ and  
  $\okC{\FF_i}{\GG_i}$, for $i=1,2$;   
\item 
$\FF_i \ArrA{\tr_i}  \FF_i'$, $i=1,2$,  with $\tr \in \interlA{\tr_1}{\tr_2}$ and 
 $\FF' = \FF'_1 \otimes \FF'_2$.
\end{enumerate} 
Then we also have 
$\GG_i \ArrA{\tr_i}   \GG_i'$, for some $\GG_i'$, 
$i=1,2$ with 
 $\GG' = \GG'_1 \otimes \GG'_2$.
\end{lemma} 


\begin{proof}
  Using Lemma~\ref{l:decNsig}, and proceeding by induction on the length of the trace.  
\end{proof} 


\begin{lemma}
\label{l:decACN}
 Suppose 
 \begin{enumerate}
   \item $\FF$ is terminating;
   \item 
   $\FF  \TEc \GG$; 
   \item 
   $\FF = \FF_1 \otimes \FF_2$,   
   and  
   $\GG = \GG_1 \otimes \GG_2$;
   \item $\okC{\FF_1}{\FF_2}$ and $\okC{\GG_1}{\GG_2}$.
 \end{enumerate} 
 Then
 $\FF_i \TEc \GG_i$, $i=1,2$. 
\end{lemma}  

\begin{proof}
  From the fact that $\FF$ is terminating and 
  $\FF \TEc \GG$, we deduce the existence of 
  a complete trace $t$ for both $\FF$ and $\GG$.
  Applying Lemma~\ref{l:AOGS-decomp-ct},
  we deduce for both $i=1,2$ the existence of a complete trace $t_i$
  for both $\FF_i$ and $\GG_i$. 
 
  Taking $j=1,2$ and $\tr_j$ a complete trace for $\FF_j$,
  from Lemma~\ref{l:AOGS-recomp-ct}
  we deduce the existence of a complete trace $\tr$ for $\FF$, with:
  \begin{itemize}
    \item either $\tr \in \interlA{\tr_1}{t_2}$ if $j=1$;
    \item or $\tr \in \interlA{t_1}{\tr_2}$ if $j=2$.
  \end{itemize}
  From $\FF  \TEc \GG$, we deduce that $\tr$ is a complete trace for $\GG$,
  so that applying Lemma~\ref{l:decNsig2},
  we get that $\tr_j$ is a complete trace for $\GG_j$.
\end{proof}
\noindent
As explained in Remark~\ref{r:CONFnoSINGL}, 
this previous lemma fails for ordinary (not necessarily complete) trace equivalence.
What makes it true for complete trace equivalence is the additional hypothesis
that $\FF$ is terminating.
%
In the other direction, a form of compositionality  
for $\TEc$ \wrt $\otimes$ also holds.

\begin{lemma}
\label{l:compCN}
Suppose 
\begin{enumerate}
\item 
$\FF = \appendCon{\FF_1}{\FF_2}$,   
 and  
$\GG = \appendCon{\GG_1}{\GG_2}$;
\item 
${\FF_i}  \TEc {\GG_i}$, $i=1,2$. 
\end{enumerate} 
Then 
${\FF} \TEc {\GG}$.
\end{lemma}              


\begin{proof}
First note that $\okC{\FF}{\GG}$ holds, 
as ${\FF_i}  \TEc {\GG_i}$, $i=1,2$ and therefore 
$\okC{{\FF_i}}{{\GG_i}}$. 

Suppose $\FF \ArrA \tr \FF'$ where $\tr$ is a complete trace for $\FF$. 
By Lemma~\ref{l:AOGS-decomp-ct}, we have 
${\FF_i} \ArrA{\tr_i} \FF_i'$,
where $ \tr_i$ is a complete trace, $i=1,2$,
with $\tr \in \interlA{\tr_1}{\tr_2} $ and  
$\FF' = \appendCon{\FF'_1}{\FF'_2}$.
%
Since ${\FF_i}  \TEc {\GG_i}$, $i=1,2$, we also have  
${\GG_i} \ArrA{\tr_i}   \GG_i'$. 
By Lemma~\ref{c:dec_GS}(\ref{l:dec_GS_2}) we infer that 
$\GG \ArrA{\tr}$. Moreover, $\tr$ is complete for $\GG$,
since $\GG$ has the same continuation structure as $\FF$.
\end{proof}  

\subsection{\WBOGS}
\label{a:tensor-wbogs}

  
We now introduce the tensor product of stacked configurations.
Compared to the tensor product of \COGS or \AOGS configurations,
there are multiple ways in \WBOGS
to tensor two stacked configurations, so that the definition of the 
tensor product $\otimes_{\stt}$ has to be 
indexed by a stack of continuation names $\stt$.
Indeed, when tensoring two stacked configurations, we can
choose how to interleave their stack.
% \begin{definition}
%   An interleaving $\stt$ of two disjoint valid stacks
%    $\stt_1,\stt_2$
%   is \emph{valid} if $\stt$ alternates between P and O continuation names
%   (\wrt a polarity function $\Label)$,
%  and for all subword $pq$ of $\stt$ with $p$ a P-name,
%  $pq$ is a subword of one of $\stt_1,\stt_2$.
% \end{definition}

First, we adapt the notion of compatible configuration
to stacked ones.
\begin{definition}
 Two stacked configurations $\FFb_1,\FFb_2$ are compatible 
 when their stacks are disjoint and 
 their corresponding \AOGS
 configurations $\es{\FFb_1},\es{\FFb_2}$ are compatible.
\end{definition}


\begin{definition}
Let $\FFb_1$ and $\FFb_2$ be two compatible stacked configurations
with at most one active,
and $\stt_1,\stt_2$ their respective stacks.  
Taking $\stt$ an interleaving of $\stt_1,\stt_2$,
we write $\FFb_1 \otimes_{\stt} \FFb_2$ for the following 
stacked configuration:
\begin{itemize}
 \item $\conf{M,p,\stt,\gamma_1 \cdot \gamma_2,\phi_1 \cup \phi_2}$
 when one of the two configurations is active, with $M,p$
 its toplevel term and continuation name;
 %and the toplevel name of $\stt$ is the toplevel one of the
 %stack of the active configuration;
 \item $\conf{\stt,\gamma_1 \cdot \gamma_2,\phi_1 \cup \phi_2}$
 when the two configurations are passive.
\end{itemize}
with $\gamma_i,\phi_i$ the environment and the set of names in $\FFb_i$.
\end{definition}
\noindent
Taking $\tr_1,\tr_2$ two traces well-bracketed 
\wrt respectively $\stt_1$ and $\stt_2$,
and $\stt$ a valid interleaving of 
$\stt_1,\stt_2$;
we write $\interlWBs{\stt}{\tr_1}{\tr_2}$ for the subset of $\interlA{\tr_1}{\tr_2}$
formed by well-bracketed traces \wrt $\stt$.
%and $\interlWB{\tr_1}{\tr_2}$ for the set
%of well-bracketed traces \wrt any valid interleaving of $\stt_1$ and $\stt_2$.


\begin{lemma}
  \label{c:lift-atrace-wbtrace}
  Suppose $\FFb_1,\FFb_2$ are compatible stacked configurations
  and $\stt$ is an interleaving of their respective stack,
  such that $\FFb_1 \otimes_{\stt} \FFb_2 \ArrWB{\tr} \GGb$.
  Taking $\tr_1,\tr_2$ two traces, $\tr \in \interlWBs{\stt}{\tr_1}{\tr_2}$,
  and $\GG_1,\GG_2$ two compatible \AOGS configurations
  we suppose that $\es{\FFb_1} \ArrA{\tr_1} \GG_1$
  and $\es{\FFb_2} \ArrA{\tr_2} \GG_2$.
% 
  Then there exists $\GGb_1,\GGb_2$ two compatible stacked configurations
  and $\stt'$ an interleaving of their respective stack such that
  $\GGb = \GGb_1 \otimes_{\stt'} \GGb_2$
  and for all $i=1,2$, 
  $\es{\GGb_i} = \GG_i$
  and 
  $\FFb_i \ArrWB{\tr_i} \GGb_i$.
\end{lemma}

\begin{proof}
We reason by induction on $\tr$.
If $\tr$ is empty, then $\GGb = \FFb_1 \otimes_{\stt} \FFb_2$.
Then both $\tr_1,\tr_2$ are empty as well,
so that $\GG_1 = \es{\FFb_1}$ and $\GG_2 = \es{\FFb_2}$.
So we define $\GGb_1$ as $\FFb_1$ and $\GGb_2$ as $\GGb_2$.

If $\tr$ is of the shape $\act \tracerho$,
then either $\act$ is the first action of $\tr_1$, or the first action of $\tr_2$.
Suppose it is the first action of $\tr_1$, then $\tr_1$
can be written as $\act \tracerho_1$.
Then there exists $\FF'_1$ such that 
$\es{\FFb_1} \arrA{\act} \FF'_1 \ArrA{\tracerho_1} \GG_1$.
Moreover, there exists $\FFb'$
such that $\FFb_1 \otimes_{\stt} \FFb_2 \arrWB{\act} \FFb' \ArrWB{\tr} \GGb$.

We prove by a case analysis over $\act$
that there exists $\FFb'_1$ and $\stt''$
such that $\FFb' = \FFb'_1 \otimes_{\stt''} \FFb_2$,
$\es{\FFb'_1} = \FF'_1$
and that $\tracerho \in \interlWBs{\stt''}{\tracerho_1}{\tr_2}$.
\begin{itemize}
  \item If $\act$ is an Opponent answer
  $\ansO{p}{x}$ then $\FFb_1$ is of the shape
  $\conf{\stt_1,\gamma\cdot [p \mapsto (E,q)],\phi}$,
  and $\stt$ is of the shape $p::\stt''$.
  So $\stt_1$ can be written as $p::\stt'_1$
  since $\stt$ is an interleaving of $\stt_1$ and $\stt_2$,
  and that $p$ must belong to $\stt_1$.
  Writing the stack of $\FFb_2$
  as $\stt_2$, we get that
  $\stt''$ is an interleaving of $\stt'_1$ and $\stt_2$.
  Defining $\FFb'_1$ as 
  $\conf{\stt'_1,\gamma,\phi \cup \{x\}}$,
  we then get the wanted properties.
  %%%%%%%%%%%%%%%%%%%%%%%%%%%
  \item If $\act$ is a Player question
  $\questPV{x}{y}{q}$
  then $\FFb_1$ is of the shape
  $\conf{E[xV],p,\stt_1,\gamma,\phi}$,
  so that $\FF'_1$ is equal to
  $\conf{\gamma \cdot \pmap{y}{V}\cdot \pmap{q}{(E,p)},\phi \cup \{y,q\}}$.
  Moreover, the stack of $\FFb'$ is equal to
  $q::\stt$.
  Defining $\FFb'_1$ as 
  $\conf{q::\stt_1,\gamma \cdot \pmap{y}{V}\cdot \pmap{q}{(E,p)},\phi \cup \{x\}}$
  and $\stt''$ as $q::\stt$,
  we then get the wanted properties.
  \item If $\act$ is an Opponent question or a Player answer,
  then we can simply take $\stt'$ equal to $\stt$,
  and $\FFb'_1$ as $\FF'_1$ equipped with the same stack
  as $\FFb_1$.
\end{itemize}
\noindent
We then conclude using the induction hypothesis over $\tracerho$.
A similar reasoning applies when $\act$ is the first action of $\tr_2$.
\end{proof}


\begin{lemma}
\label{c:dec_GS_WB}
Suppose $\FFb_1,\FFb_2$ are compatible stacked configurations
and $\stt$ is an interleaving of their respective stack.
Then the set of traces generated by $\FFb_1 \otimes_{\stt} \FFb_2$
is equal to the set of interleavings 
$\interlWBs{\stt}{\tr_1}{\tr_2}$ well-bracketed \wrt $\stt$,
with $\FFb_1 \ArrWB{\tr_1}$ and $\FFb_2 \ArrWB{\tr_2}$.
Moreover, $\tr$ is complete if and only if $\tr_1,\tr_2$ are.
That is:
\begin{enumerate}
  \item If $\FFb_1 \otimes_{\stt} \FFb_2 \ArrWB \tr  \GGb$, then, 
  for $i=1,2$,  there are traces $\tr_i$
  such that ${\FFb_i}\ArrWB{\tr_i} \GGb_i$, 
  and $\tr \in \interlWBs{\stt}{\tr_1}{\tr_2}$, 
  with $\GGb = \GGb_1 \otimes_{\stt'} \GGb_2$;
  and if $\tr$ is complete for $\FFb$ then both $\tr_i$ are complete for $\FFb_i$;
  %
  \item Conversely, if for  $i=1,2$, we have
  ${\FFb_i} \ArrWB {\tr_i} \GGb_i$,  
  and $\tr \in \interlWBs{\stt}{\tr_1}{\tr_2}$, 
  then $\FFb_1 \otimes_{\stt} \FFb_2 \ArrWB{\tr} 
  \GGb_1 \otimes_{\stt'} \GGb_2$, with $\stt'$
  an interleaving of the stack of $\GGb_1$ and $\GGb_2$;
  and if both $\tr_i$ are complete for $\FFb_i$ then $\tr$ is complete for $\FFb$.
\end{enumerate} 
\end{lemma}

\begin{proof}
  \begin{enumerate}
    \item Lifting the reduction to \AOGS we get that
    $\es{\FFb_1} \otimes_{\stt} \es{\FFb_2} \ArrA{\tr} \es{\GGb}$.
    Applying Lemma~\ref{c:dec_GS} we deduce that there exists
    $\GG_1,\GG_2$ two compatible \AOGS configurations
    such that $\es{\FFb_1} \ArrA{\tr_1} \GG_1$
    and $\es{\FFb_2} \ArrA{\tr_2} \GG_2$.
    We conclude using Lemma~\ref{c:lift-atrace-wbtrace}.
    \item Lifting the reductions to \AOGS we get that
    ${\es{\FFb_i}} \ArrA {\tr_i} \es{\GGb_i}$, $i=1,2$.
    Applying Lemma~\ref{c:dec_GS} we deduce that 
    $\es{\FFb_1} \otimes \es{\FFb_2} \ArrA{\tr} \es{\GGb_1} \otimes \es{\GGb_2}$.
    Since $\tr$ is well-bracketed,
    from Lemma~\ref{l:WBs_lts}
    we get that
    $\FFb_1 \otimes_{\stt} \FFb_2 \ArrWB{\tr} 
     \GGb_1 \otimes_{\stt'} \GGb_2$
    for $\stt'$
    an interleaving of the stack of $\GGb_1$ and $\GGb_2$.

\noindent
    If both $\tr_i$ are complete, then from Lemma~\ref{l:WBOGS_complete-strongly-passive},
    both $\GGb_i$ are strongly passive, so that $\GGb$ is also strongly passive.
    Thus applying Lemma~\ref{l:WBOGS_complete-strongly-passive} in the other direction, 
    we deduce that $\tr$ is complete for $\FFb$.
    \qedhere
  \end{enumerate}
\end{proof}
\noindent
Finally, we adapt Lemma~\ref{l:decACN} to \WBOGS.


\ifcutLMCS
Lemma~\ref{l:decNsig-WB} shows that the action of a configuration is 
determined by a specific singleton component, 
and it is the action itself that allows us to determine such a component.

\begin{lemma}
\label{l:decNsig-WB}
Suppose $\FFb \arrWB\mu \FFb'$.
Then there is a singleton $\FFb_1$, and some $\FFb_1',\FFb_2,\FFb_2'$ such that:

\begin{enumerate}
\item $\FFb =
 {\FFb_1} \otimes_{\stt} {\FFb_2}$,
\item $\FFb_1 \arrWB\mu \FFb_1'$,
\item $\FFb' =
 {\FFb_1'} \otimes_{\stt'} {\FFb_2}$.
\end{enumerate} 
 
Moreover, the P-\Supports and the structure of continuations of $\FFb_1$ and $\FFb_1'$  
are determined by 
$\mu$; that is, 
if  $\okC \FFb\GGb$ and
 $\GGb \arrWB\mu \GGb'$, then for  
 the singleton $\GGb_1$, and the configurations $\GGb_1',\GGb_2,\GGb_2'$ such that
 $\GGb = \appendCon
 {\GGb_1} {\GGb_2}$,
 $\GGb_1 \arrWB\mu \GGb_1'$, and 
 $\GGb' = \appendCon
 {\GGb_1'} {\GGb_2}$ we have: 
 $\okC{\FFb_1}{\GGb_1}$, $\okC{\FFb_2}{\GGb_2}$, and 
$\okC{\FFb_1'}{\GGb_1'}$.  
 \end{lemma} 



\begin{proof}
The first part is proved using Lemma~\ref{c:dec_GS_WB}. 
The second part is a simple case analysis on the form of $\mu$.
\end{proof} 

The next lemma uses the previous result to show that the decomposition of a trace 
is uniform.

\begin{lemma}
\label{l:decNsig2-WB}
Suppose 
\begin{enumerate}
\item
$\FFb \ArrWB \tr   \FFb'$
and $\GGb \ArrWB \tr   \GGb'$, 
\item  
 $\FFb = \FFb_1 \otimes_\stt \FFb_2$ and 
 $\GGb = \GGb_1 \otimes_\stt \GGb_2$,  
\item 
  $\okC \FFb\GGb$ and  
  $\okC {\FFb_i}{\GGb_i}$, for $i=1,2$,   
\item 
$\FFb_i \ArrWB{\tr_i}  \FFb_i'$, $i=1,2$,  with $\tr \in \interlWBs{\stt}{\tr_1}{\tr_2}$ and 
 $\FFb' = \FFb'_1 \otimes_{\stt'} \FFb'_2$.
\end{enumerate} 
Then we also have 
$\GGb_i \ArrWB{\tr_i}   \GGb_i'$, for some $\GGb_i'$, 
$i=1,2$ with 
 $\GGb' = \GGb'_1 \otimes_\stt \GGb'_2$.
\end{lemma} 


\begin{proof}
Using Lemma~\ref{l:decNsig-WB}, and proceeding by induction on the length of the trace.  
\end{proof} 
\fi

\begin{lemma}
\label{l:decWBCN}
  Suppose 
  \begin{enumerate}
    \item $\FFb$ is terminating;
    \item $\FFb  \TEwbc \GGb$; 
    \item $\FFb = \FFb_1 \otimes_{\stt} \FFb_2$,   
    and  
    $\GGb = \GGb_1 \otimes_{\stt} \GGb_2$;
    \item $\okC{\FFb_1}{\FFb_2}$ and $\okC{\GGb_1}{\GGb_2}$.
  \end{enumerate} 
  Then
  $\FFb_i  \TEwbc \GGb_i$, $i=1,2$. 
\end{lemma}   
\noindent
Its proof is a direct adaption of the proof of Lemma~\ref{l:decACN},
using Lemmas corresponding to Lemma~\ref{l:decNsig} and Lemma~\ref{l:decNsig2}
for \WBOGS.


\ifcutLMCS

\begin{proof}
  We write $\stt_1$ for the stack of both $\FFb_1, \GGb_1$
  and $\stt_2$ for the one of $\FFb_2, \GGb_2$.
  Suppose $\FFb_1  \TEwbc \GGb_1$ 
  does not hold (the  case of $\FFb_2$ and $\GGb_2$  is similar).
  Thus there is $\tracerho$ a well-bracketed complete trace \wrt $\stt_1$, for $\FFb_1$ and not for $\GGb_1$. 
  
  Since $\FFb$ is terminating
  and $\FFb  \TEwbc \GGb$, 
  there is for both $\FFb$ and $\GGb$ 
  a well-bracketed complete trace
  $\tr$ \wrt $\stt$
  the stack of $\FFb$ and $\GGb$.
  
  By Lemma~\ref{c:dec_GS_WB}, there are traces $\tr_i$
  well-bracketed complete \wrt $\stt_i$ such that 
  $\FFb_i \ArrWB {\tr_i}$ 
  and $\tr \in \interlWBs{\stt}{ \tr_1}{\tr_2}$.
  %
  By Lemma~\ref{l:decNsig2-WB}, we also have that
  $\GGb_i \ArrWB {\tr_i}$.   
  
  By Lemma~\ref{c:dec_GS_WB}, there is 
   $\tracetheta \in \interlWBs{\stt}{\tracerho}{\tr_2}$ complete such that
  $\FFb \ArrWB {\tracetheta } $.
  Since 
  $\FFb  \TEwbc \GGb$, 
  we also have $\GGb \ArrWB {\tracetheta } $. 
  By Lemma~\ref{l:decNsig2-WB}, 
  we deduce that 
  $\GGb_1 \ArrWB {\tracerho }$ and 
  $\GGb_2 \ArrWB {\tr_2 }$.
  This contradicts the assumption that 
   ${\tracerho }$ is not a trace for $\GGb_1$.
\end{proof} 
\fi